\documentclass[10pt,a4paper]{amsart}
\usepackage[margin=19mm]{geometry}
\usepackage{amsmath,amssymb,mathtools}
\usepackage[scr=rsfs,cal=euler]{mathalfa}
\usepackage{xfrac}
\usepackage[unicode,hidelinks,bookmarks=false]{hyperref}
\newcommand{\ie}{i.e.\ }
\newcommand{\cf}{cf.\ }
\newcommand{\resp}{respectively\ }
\newcommand{\Subsection}{\subsection}
\providecommand{\nobreakdash}{\nobreak\hbox{-}}
\setlength{\emergencystretch}{2em}
\multlinegap0pt
\usepackage{bm}
\usepackage{bbm}
\usepackage{dsfont}
\usepackage{xspace}
\usepackage{cancel}
\usepackage{mathtools}
\usepackage{amsthm}
\makeatletter
\@ifundefined{lemma}{\newtheorem{lemma}{Lemma}}{}
\makeatother
\title[Fractional p-Laplacian Kirchhoff-type problem involving a si]{Fractional p-Laplacian Kirchhoff-type problem involving a singular term via Nehari manifold}
\author{Djamel Abid}
\date{Source: arXiv:2510.07911v1 (CC BY 4.0)}
\begin{document}
\maketitle

\noindent\emph{Source and licence.} Fractional p-Laplacian Kirchhoff-type problem involving a singular term via Nehari manifold. Version arXiv:2510.07911v1 (CC BY 4.0), distributed under the Creative Commons Attribution 4.0 International licence, as recorded in the arXiv licence metadata for this exact version and shown on the abstract page https://arxiv.org/abs/2510.07911v1. Full rights notice: licenses/arxiv-2510.07911-CC-BY-4.0.txt.\par\smallskip

\noindent\textit{Editorial scope.} Counted unit: the Lemma whose source label is \texttt{l7} in the article (source0.tex lines 660--663, source label \texttt{l7}), delivered together with the article's own complete original proof of it (source0.tex lines 664--778, one \texttt{proof} environment of 115 lines). No mathematics is written, repaired or re-derived: statement and proof are the article's own text line for line. Required context retained uncounted: l4 (lines 468--475); ek2 (lines 585--587); l6 (lines 601--613). Every definition, display and label of those lines is retained unchanged. Omitted from this package and not redistributed: 1--467, 476--584, 588--600, 614--659, 779--1530. Nothing inside a retained excerpt is omitted or altered: every retained body is delivered exactly as the article's lines are written, so no omission receipt of any kind accompanies this package. Citations: the retained excerpts carry no citation command at all, so no bibliography entry of the article is transcribed and no reference is redistributed. Reference adaptation: none was needed and none was made; every reference inside the delivered unit resolves inside it, so no unresolved reference or citation is produced. Printed slips kept literal: no printed word, symbol, hypothesis or number of the counted statement or of its proof was altered. Layout and class: the article's own class is \texttt{amsart} with packages including color, amsmath, amsthm, amsfonts, amssymb, breqn; the wrapper is the lane's pinned \texttt{amsart} editorial wrapper at 10pt on A4, which restates the theorem-like environments the retained text needs and adds the article's own macro definitions that the retained text needs (none), with extra wrapper packages (bm, bbm, dsfont, xspace, cancel, mathtools). The delivered numbering is the unit's own. Assets: no retained span contains a figure environment, a graphics-inclusion command, a PDF-page inclusion, a table or a TikZ picture; no image file is copied into this package, the package declares no asset dependency and no image file is redistributed with it. Licence: version arXiv:2510.07911v1 of this article is distributed under the Creative Commons Attribution 4.0 International licence, as recorded in the arXiv licence metadata for this exact version and shown on the abstract page https://arxiv.org/abs/2510.07911v1. A full rights notice is delivered at licenses/arxiv-2510.07911-CC-BY-4.0.txt. Characters: every retained excerpt is pure ASCII, so the delivered engine, XeLaTeX with its default Latin Modern text font, reports no missing character.
\begin{lemma}\label{l4}
Let $ u \in \mathcal{N}_{\lambda}^{-},$ and $ \lambda > 0.$ Then there exist $ \epsilon > 0$ and a continuous function 
$ \zeta : B_{\epsilon}(0) \longrightarrow \mathbb{R}^{+}$ such that
\begin{equation*}
\zeta(\phi) > 0, \zeta(0) = 1, \,\,\, \zeta(\phi)(u +\phi)\,\,\, in\,\,\, \mathcal{N}_{\lambda}^{-} \,\,\, \textit{for all} \,\,\, \phi \,\,\, \in \,\,\, B_{\epsilon}(0).
\end{equation*}
Where  $ B_{\epsilon}$ is the open ball in $ X_0$ centered at $ 0$ and of radius $ \epsilon.$
\end{lemma}

\begin{equation}\label{ek2}
J_{\lambda}(u) \geq J_{\lambda}(u_n) -\frac{1}{n} \Vert u -u_n \Vert, \quad \textit{for all} \quad u \in \mathcal{N}_{\lambda},
\end{equation}

\begin{lemma}\label{l6}
Let $ \lambda \in (0, \lambda_*), (u_n)_n \in \mathcal{N}$ we have the following inequalities hold for any $ n \in \mathbb{N},$
\begin{itemize}
\item If $ (u_n)_n \in \mathcal{N}_{\lambda}^{+},$ then
\begin{equation*}
(p +\alpha -1) \left(a\Vert u_n \Vert^p + \Vert u_n \Vert_{p}^{p}\right) + b(p \theta +\alpha -1)\Vert u_n \Vert^{p \theta} - \lambda q (q +\alpha -1) \int F(x,u_n^+) >0.
\end{equation*}
\item If $ (u_n)_n \in \mathcal{N}_{\lambda}^{-},$ then
\begin{equation*}
(p +\alpha -1) \left(a\Vert u_n \Vert^p + \Vert u_n \Vert_{p}^{p}\right) + b(p \theta +\alpha -1)\Vert u_n \Vert^{p \theta} - \lambda q (q +\alpha -1) \int F(x,u_n^+) <0.
\end{equation*}
\end{itemize}
\end{lemma}

\begin{lemma}\label{l7}
Let $ (u_n)_n \in \mathcal{N}_{\lambda}^{-},$ (resp $ \mathcal{N}_{\lambda}^{+},$) such that up to subsequence  $ u_n \rightharpoonup u_*$ on $ X_0$ and satidfying 
\ref{ek2}, then for $ \lambda \in (0, \lambda_*)$ we have, $ \langle\zeta^{\prime}(0), \Phi \rangle$ is uniformly bounded for any positive $ \Phi$ in $ X_0.$
\end{lemma}

\begin{proof}
Let $ u_n \in \mathcal{N}_{\lambda}^{+},$ then we have
\begin{equation}\label{e08}
\left( a \Vert u_n \Vert^p + \Vert u_n \Vert_{p}^{p}\right) + b \Vert u_n \Vert^{p \theta}  - \int c(x)(u_n^+)^{-\alpha+1} - \lambda q \int F(x,u_n^+) = 0 
\end{equation}
and by lemma \ref{l4}, for $ \Phi \in X_0,$ where $ \Phi \geq 0,$ there exist a sequence of functions $ (\zeta_n)_n$ such that 
$ \zeta_n(t\Phi(u_n +t\Phi)) \in \mathcal{N}_{\lambda}^{+}$ and $ \zeta_n(0) =1.$ which yields
\begin{equation}\label{e09}
\begin{split}
\zeta_n^{p}(t\Phi) \left( a \Vert u_n +t\Phi \Vert^p + \Vert u_n +t\Phi \Vert_{p}^{p}\right) + b \zeta_n^{p\theta}(t\Phi) \Vert u_n +t\Phi \Vert^{p \theta}\\ 
-\zeta_n^{-\alpha+1}(t\Phi) \int c(x)((u_n +t\Phi)^+)^{-\alpha+1} - \lambda \zeta_n^{q}(t\Phi) q \int F(x,(u_n +t\Phi)^+) = 0 
\end{split}
\end{equation} 
Where $ \langle\zeta_n^{\prime}(0), \Phi\rangle \in \bar{\mathbb{R}},$ for any $ \Phi \in X_0,$ assuming that the right derivative of $ \zeta$ at $ t =0$ exists,
if not we take $ t_n >0$ such that $ t_n \xrightarrow{>} 0,$ and then we take $ \zeta_n^{\prime}(0)$ as $ \lim_{n \to \infty} s_n =\frac{\zeta_n(t_n\Phi) -1}{t_n}.$ Therefore by \ref{e08} and \ref{e09} we obtain
\begin{equation}
\begin{aligned}
&(\zeta_n^{p}(t\Phi) -1) \left(a\Vert u_n +t\Phi \Vert^p + \Vert u_n +t\Phi \Vert_{p}^{p}\right)+a\left(\Vert u_n +t\Phi \Vert^p -\Vert u_n \Vert^p\right)\\ 
& +\left(\Vert u_n +t\Phi \Vert_{p}^{p}-\Vert u_n \Vert_{p}^{p}\right) +b (\zeta_n^{p\theta}(t\Phi) -1) \Vert u_n +t\Phi \Vert^{p \theta}  +b\left(\Vert u_n +t\Phi \Vert^{p \theta} -\Vert u_n \Vert^{p \theta}\right)\\
&-(\zeta_n^{-\alpha+1}(t\Phi) -1)\int c(x)((u_n +t\Phi)^+)^{-\alpha+1}\\
&-\left(\int c(x)((u_n +t\Phi)^+)^{-\alpha+1} -\int c(x)\vert u_n\vert^{-\alpha+1}\right)\\
&-\lambda q (\zeta_n^{q}(t\Phi) -1) \int F(x,(u_n +t\Phi)^+)\\
& -\lambda q \left(\int F(x,(u_n +t\Phi)^+) -\int F(x,u_n^+)\right) =0\\
\end{aligned}
\end{equation} 
deviding by $ t >0$ and passing to the limit when $ t \xrightarrow{>} 0,$ and usign \ref{e08} we find,
\begin{equation}
\begin{aligned}
& 0\leq p\zeta^{p-1}_{n}(t\Phi)\langle \zeta^{\prime}_{n}(0), \Phi\rangle\left( a \Vert u_n \Vert^p + \Vert u_n \Vert_{p}^{p}\right)\\
&+pa\int_{\Omega}\frac{\vert u_n(x) -u_n(y)\vert^{p-2}(u_n(x) -u_n(y))(\Phi(x) -\Phi(y))}{\vert x -y\vert^{n  +ps}}dxdy \\
&+\int_{\Omega}(u_n^+)^{p -1}\Phi dx +b p\theta\zeta^{p\theta-1}_{n}(t\Phi)\langle \zeta^{\prime}_{n}(0), \Phi\rangle\Vert u_n \Vert^{p\theta} 
+b p\theta\Vert u_n\Vert^{p(\theta -1)}\\
&\int_{\Omega}\frac{\vert u_n(x) -u_n(y)\vert^{p-2}(u_n(x) -u_n(y))(\Phi(x) -\Phi(y))}{\vert x -y\vert^{n  +ps}}dxdy\\
&-(-\alpha +1)\zeta^{-\alpha}_{n}(t\Phi)\langle \zeta^{\prime}_{n}(0), \Phi\rangle\int c(x)(u_n^+)^{-\alpha+1}\\
& -(-\alpha +1)\int_{\Omega}c(x)(u_n^+)^{-\alpha}\Phi -\lambda q^{2}\zeta^{q-1}_{n}(t\Phi)\langle \zeta^{\prime}_{n}(0), \Phi\rangle\int_{\Omega}F(x,u_n^+)dx\\
&-\lambda q\int_{\Omega}f(x,u_n^+)\Phi dx\\
&= \langle \zeta^{\prime}_{n}(0), \Phi\rangle\\
&\left( (p +\alpha -1)\left( a \Vert u_n \Vert^p + \Vert u_n \Vert_{p}^{p}\right) +b(p\theta +\alpha -1) -\lambda q(q +\alpha -1)\int F(x,u_n^+)dx\right)\\
& +\left(a p +b p\theta\Vert u_n\Vert^{p(\theta -1)}\right)\int_{\Omega}\frac{\vert u_n(x) -u_n(y)\vert^{p-2}(u_n(x) -u_n(y))(\Phi(x) -\Phi(y))}{\vert x -y\vert^{n  +ps}}dxdy \\
&-\lambda q\int_{\Omega}f(x,u_n^+)\Phi dx +\int_{\Omega}(u_n^+)^{p -1}\Phi dx
\end{aligned}
\end{equation}
Since $ u_n$ is bounded in $ X_0$ and by lemma \ref{l6} there exist a constant 
$ C_1 >0$ and $ C_2 \in \mathbb{R}$ such that $ \langle \zeta^{\prime}_{n}(0) \Phi\rangle \geq \frac{C_2}{C_1},$ 
hence $ \langle \zeta^{\prime}_{n}(0) \Phi\rangle$ is bounded from bellow for any positive $ \Phi$ in $ X_0.$\newline
On the other hand using \ref{ek2} with $ u =\zeta_n(t\Phi)(u_n +t\Phi),$ we obtain
\begin{equation}
\begin{aligned}
& \vert (\zeta_n(t\Phi) -1)\frac{\Vert u_n\Vert}{n} +\zeta_n(t\Phi)\frac{\Vert t\Phi\Vert}{n} 
\geq J_{\lambda}(u_n) -J_{\lambda}(\zeta_n(t\Phi)(u_n +t\Phi))\\
& = \left( \frac{1}{p} -\frac{1}{-\alpha +1}\right)\left( a \Vert u_n \Vert^p + \Vert u_n \Vert_{p}^{p}\right)\\
& +b\left( \frac{1}{p\theta} -\frac{1}{-\alpha +1}\right)\Vert u_n \Vert^{p \theta}\\
& +\lambda\left( \frac{q +\alpha -1}{-\alpha +1}\right) \int F(x,u_n^+)\\
& +\left(\frac{1}{-\alpha +1} - \frac{1}{p}\right)\zeta_n^{p}(t\Phi) \left( a \Vert u_n +t\Phi \Vert^p + \Vert u_n +t\Phi \Vert_{p}^{p}\right)\\
& +b\left( \frac{1}{-\alpha +1} -\frac{1}{p\theta}\right)\zeta_n^{p\theta}(t\Phi) \Vert u_n +t\Phi \Vert^{p \theta}\\
& -\lambda \left(  \frac{q +\alpha -1}{-\alpha +1}\right)\zeta^{q}_n(t\Phi) \int F(x,(u_n +t\Phi)^+)dx\\
& =\left(\frac{1}{-\alpha +1} - \frac{1}{p}\right)\left(a (\Vert u_n +t\Phi \Vert^p - \Vert u_n \Vert^p) +\Vert u_n +t\Phi \Vert_{p}^{p} -\Vert u_n \Vert_{p}^{p}\right)\\
& -b\left( \frac{1}{p\theta} -\frac{1}{-\alpha +1}\right)\left(  \Vert u_n +t\Phi \Vert^{p \theta} -\Vert u_n \Vert^{p \theta}\right)\\
& +\left(\frac{1}{-\alpha +1} - \frac{1}{p}\right)(\zeta_n^{p}(t\Phi) -1)\left( a \Vert u_n +t\Phi \Vert^p + \Vert u_n +t\Phi \Vert_{p}^{p}\right)\\
& +b\left( \frac{1}{-\alpha +1} -\frac{1}{p\theta}\right)(\zeta_n^{p\theta}(t\Phi) -1)\Vert u_n +t\Phi \Vert^{p \theta}\\
& -\lambda\left(  \frac{q +\alpha -1}{-\alpha +1}\right)\left( \int F(x,(u_n +t\Phi)^+)dx -\int F(x,u_n^+)dx\right)\\
& -\lambda\left(  \frac{q +\alpha -1}{-\alpha +1}\right)(\zeta^{q}_n(t\Phi) -1)\int F(x,(u_n +t\Phi)^+)dx. \\
\end{aligned}
\end{equation}

deviding by $ t$ and passing to the limmit where $ t \xrightarrow{>} 0,$ we find
{\scriptsize
\begin{equation}
\begin{aligned}
& \langle\zeta^{\prime}_n(0),\Phi \rangle\frac{\Vert u_n\Vert}{n} +\frac{\Vert \Phi\Vert}{n} \geq \left(\frac{p +\alpha -1}{p(-\alpha +1)}\right)\\
&\left(p a\int_{\Omega}\frac{\vert u_n(x) -u_n(y)\vert^{p-2}(u_n(x) -u_n(y))(\Phi(x) -\Phi(y))}{\vert x -y\vert^{n  +ps}}dxdy +\int_{\Omega}(u_n^+)^{p -1}\Phi dx\right)\\
& +\left(\frac{p +\alpha -1}{-\alpha +1}\right)\langle \zeta^{\prime}_n(0),\Phi \rangle\left( a\Vert u_n\Vert^{p} +\Vert u_n\Vert^{p}_{p}\right)\\
& +b\left( \frac{p\theta +\alpha -1}{-\alpha +1}\right)\Vert u_n\Vert^{p(\theta -1)}\int_{\Omega}\frac{\vert u_n(x) -u_n(y)\vert^{p-2}( u_n(x) -u_n(y))(\Phi(x) -\Phi(y))}{\vert x -y\vert^{n  +ps}}dxdy\\
&+b\left( \frac{p\theta +\alpha -1}{-\alpha +1}\right)\langle \zeta^{\prime}_n(0),\Phi \rangle\Vert u\Vert^{p\theta }\\
& -\lambda\left(  \frac{q +\alpha -1}{-\alpha +1}\right)\int_{\Omega} f(x,u_n^+)\Phi dx\\
& -\lambda\left(  \frac{q +\alpha -1}{-\alpha +1}\right) q \langle \zeta^{\prime}_n(0),\Phi \rangle\int F(x,u_n^+)dx. \\
& =\frac{\langle \zeta^{\prime}_n(0),\Phi \rangle}{-\alpha +1}\\
&\left( (p +\alpha -1)\left( a\Vert u_n\Vert^{p} +\Vert u_n\Vert^{p}_{p}\right) +(p\theta +\alpha -1)\Vert u_n\Vert^{p\theta} -\lambda q(q +\alpha -1)\int F(x,u_n^+)dx.\right)\\
& +a\left(\frac{p +\alpha -1}{-\alpha +1}\right)\int_{\Omega}\frac{\vert u_n(x) -u_n(y)\vert^{p-2}(u_n(x) -u_n(y))(\Phi(x) -\Phi(y))}{\vert x -y\vert^{n  +ps}}dxdy
\\
& +\left(\frac{p +\alpha -1}{p(-\alpha +1)}\right)\int_{\Omega}(u_n^+)^{p -1}\Phi dx\\
& +b\left( \frac{p\theta +\alpha -1}{-\alpha +1}\right)\Vert u\Vert^{p(\theta -1)}\int_{\Omega}\frac{\vert u_n(x) -u_n(y)\vert^{p-2}(u_n(x) -u_n(y))(\Phi(x) -\Phi(y))}{\vert x -y\vert^{n  +ps}}dxdy\\
& -\lambda\left(  \frac{q +\alpha -1}{-\alpha +1}\right)\int_{\Omega} f(x,u_n^+)\Phi dx\\
\end{aligned}
\end{equation}
}
which means that
{\scriptsize
\begin{equation}
\begin{aligned}
&\frac{\Vert \Phi\Vert}{n} \geq \frac{\langle\zeta^{\prime}_n(0),\Phi \rangle}{-\alpha +1}\\
&\left( (p +\alpha -1)\left( a\Vert u_n\Vert^{p} +\Vert u_n\Vert^{p}_{p}\right) +(p\theta +\alpha -1)\Vert u\Vert^{p\theta} -\lambda q(q +\alpha -1)\int F(x,u_n^+)dx -\frac{(-\alpha +1)\Vert u_n\Vert}{n} \right)\\
& +a\left(\frac{p +\alpha -1}{-\alpha +1}\right)\int_{\Omega}\frac{\vert u_n(x) -u_n(y)\vert^{p-2}(u_n(x) -u_n(y))(\Phi(x) -\Phi(y))}{\vert x -y\vert^{n  +ps}}dxdy
\\
& +\left(\frac{p +\alpha -1}{p(-\alpha +1)}\right)\int_{\Omega}(u_n^+)^{p -1}\Phi dx\\
& +b\left( \frac{p\theta +\alpha -1}{-\alpha +1}\right)\Vert u_n\Vert^{p(\theta -1)}\int_{\Omega}\frac{\vert u_n(x) -u_n(y)\vert^{p-2}(u_n(x) -u_n(y))(\Phi(x) -\Phi(y))}{\vert x -y\vert^{n  +ps}}dxdy\\
& -\lambda\left(  \frac{q +\alpha -1}{-\alpha +1}\right)\int_{\Omega} f(x,u_n^+)\Phi dx\\
\end{aligned}
\end{equation}
}
By the boundedness of $ u_n$ in $ X_0,$ and lemma \ref{l6} we can choose $ n$ large enough such that
{\scriptsize
\begin{equation}\label{e010}
\begin{aligned}
&(p +\alpha -1)\left( a\Vert u_n\Vert^{p} +\Vert u_n\Vert^{p}_{p}\right) +(p\theta +\alpha -1)\Vert u\Vert^{p\theta} -\lambda q(q +\alpha -1)\int F(x,u_n^+)dx -\frac{(-\alpha +1)\Vert u_n\Vert}{n}\\
& = C_1 - \frac{(-\alpha +1)C_3}{n} >0.
\end{aligned}
\end{equation}
}
That is to say that there exist $ C_4 \in \mathbb{R}$  such that 
\begin{equation*}
\langle\zeta^{\prime}_n(0),\Phi \rangle \leq (-\alpha +1)\left( \frac{\Vert \Phi\Vert}{n} +C_4\right)\left( C_1 - \frac{(-\alpha +1)C_3}{n}\right)^{-1}.
\end{equation*}
Hence by \ref{e010} we get that  $ \langle\zeta^{\prime}_n(0),\Phi \rangle$ is bounded from above for $ \Phi\in X_0$ with $ \Phi \geq 0$ and $ n$ Large enough.
\end{proof}

\end{document}
